\section{Sonic loci}

For a regular point of the wave manifold, write
\[
U_- = \bar U-\frac{Y}{2}r(z),
\qquad
U_+ = \bar U+\frac{Y}{2}r(z),
\qquad
r(z)=\binom{z}{1},
\]
and let $s$ be the Rankine--Hugoniot speed.  The left and right sonic
conditions are defined intrinsically by
\[
\boxed{\det\!\left(DF(U_-)-sI\right)=0},
\qquad
\boxed{\det\!\left(DF(U_+)-sI\right)=0}.
\]
Because the flux is quadratic, its Jacobian is affine.  Define
\[
A=DF(\bar U)-sI
\]
and
\[
M(z)=D\!\left(DF(U)r(z)\right)
 =
\begin{pmatrix}
2a_{20}z+a_{11} & a_{11}z+2a_{02}\\
2b_{20}z+b_{11} & b_{11}z+2b_{02}
\end{pmatrix}.
\]
Then
\[
DF(U_-)-sI=A-\frac{Y}{2}M,
\qquad
DF(U_+)-sI=A+\frac{Y}{2}M.
\]
By construction of the wave-manifold fiber, $r(z)$ is an eigenvector of
$DF(\bar U)$ associated with $s$, hence
\[
\det A=0.
\]
For
\[
K=A_{11}M_{22}+M_{11}A_{22}-A_{12}M_{21}-M_{12}A_{21},
\qquad D_M=\det M,
\]
the two endpoint determinants factor exactly as
\[
\boxed{
\det(DF(U_-)-sI)=\frac{Y}{4}\bigl(YD_M-2K\bigr)
},
\]
\[
\boxed{
\det(DF(U_+)-sI)=\frac{Y}{4}\bigl(YD_M+2K\bigr)
}.
\]

Therefore the characteristic surface
\[
\mathcal C=\{Y=0\}
\]
is the common zero-jump factor in both endpoint determinant equations.
The sonic surfaces $\mathcal S^-$ and $\mathcal S^+$ used by the application
are the nontrivial factors:
\[
\boxed{\mathcal S^-:\;YD_M-2K=0},
\qquad
\boxed{\mathcal S^+:\;YD_M+2K=0}.
\]
Whenever $D_M\neq0$ they are graphs over $(\tau,z)$,
\[
Y_-=\frac{2K}{D_M},
\qquad
Y_+=-\frac{2K}{D_M},
\]
so that $Y_++Y_-=0$.

If $D_M=0$, division by $D_M$ is not valid.  The implementation therefore
keeps the implicit equations as the authoritative definition.  In this
case, if $K\neq0$ there is no finite point of the corresponding nontrivial
sheet at that $(\tau,z)$; if $K=0$ the factor is degenerate and requires
separate local analysis.

The coincidence locus $\mathcal E$ remains a subset of $\mathcal C$; no
additional inclusion between $\mathcal E$ and the nontrivial sheets
$\mathcal S^\pm$ is assumed here.

\subsection{Implicit geometry, inflection and double sonic locus}
The nontrivial sonic sheets are retained in implicit form
\[
 \Phi_-(\tau,Y,z)=Y\det M-2K=0,\qquad
 \Phi_+(\tau,Y,z)=Y\det M+2K=0.
\]
The formulas $Y_-=2K/\det M$ and $Y_+=-2K/\det M$ are valid only when
$\det M\neq0$ and are not used to define or render the global sonic geometry.
Indeed, simultaneous membership in both nontrivial sonic sheets gives
\[
 K=0,\qquad Y\det M=0.
\]
On the characteristic surface $Y=0$, the common locus is the inflection locus
\[
 \mathcal J=\{Y=0,\ K=0\}.
\]
Away from the characteristic surface, $Y\neq0$, the common locus is the double
sonic set
\[
 \mathcal{DS}=\{K=0,\ \det M=0,\ Y\neq0\}.
\]
At $\mathcal{DS}$ the explicit quotient is the indeterminate form $0/0$; this
is a degeneracy of that parametrization, not a disappearance of the sonic set.
